Identity testing is a fundamental algorithmic question with numerous applications.  In \emph{algebraic} identity testing, the task is to determine the zeroness of an expression evaluated in a given ring.  This problem has many different versions, depending on the syntax for giving the expression and the ring in which the evaluation is to be performed.

A basic instance of algebraic identity testing is the Arithmetic Circuit Identity Testing (ACIT) problem, which involves deciding the zeroness of an integer represented as an arithmetic circuit.  The difficulty in this problem is that the integer may have bit-length exponential in the size of the circuit.  However, the problem admits a randomised polynomial-time algorithm: one evaluates the circuit modulo a prime that is randomly chosen in a certain range.  The ACIT problem turns out to be polynomial-time interreducible with the problem of determining zeroness of an arithmetic circuit evaluated in the  ring of multivariate polynomials: the so-called Polynomial Identity Testing (PIT) problem~\cite{Allender06onthe}.  
Over the years, {PIT} has found diverse applications in algorithm design; a few well-known examples are program testing~\cite{demillo-lipton}, detecting perfect matchings~\cite{lovasz}, factoring polynomials~\cite{kaltofen}, pattern matching in compressed texts~\cite{konig-lohrey, berman-pattern}, primality testing~\cite{agrawal-biswas, aks}, equivalence and minimization of weighted automata~\cite{kiefer-minimize-wa, kiefer-marusic-worrell} and linear recurrence sequences~\cite{COW,abmvv}.
Whether PIT admits a deterministic polynomial-time algorithm is one of the central open questions in complexity theory.

The topic of this paper is \emph{radical identity testing}, that is, testing zeroness of an expression in radicals, represented by an algebraic circuit.
This generalizes the ACIT problem: the evaluation of the circuit occurs in the ring of integers of a number field, rather than the ring of integers of the rational numbers.  Historically, expressions in radicals played a central role in solving polynomial equations.  Such expressions also naturally arise in optimization problems on graphs embedded in Euclidean space, such as the Euclidean Minimal Spanning Tree and Traveling Salesperson problems \cite{ggj}.  Another source of radical expressions arises from an approach by Chen and Kao~\cite{chen-kao} to derandomizing PIT.  Their idea was to test the zeroness of a multivariate polynomial by evaluating it on a certain randomly chosen radical expression.  In their method, by construction, the radical expressions are such that the result of the evaluation is zero if and only if the polynomial is identically zero.  The challenge then becomes to determine the zeroness of the resulting expression in radicals, for which Chen and Kao use numerical approximation.  This approach allowed to reduce the number of random bits required for the problem compared with previously known methods, such as those based on the Schwarz-Zippel lemma.

In this paper, we introduce a symbolic approach to the problem of identity testing algebraic integers in the number field generated over $\mathbb{Q}$ by \emph{real radicals}; given an algebraic circuit  representing a multivariate polynomial $f(x_1,\ldots,x_k)\in\ZZ[x_1,\ldots,x_k]$, and  radical inputs $\mysqrt{d_1}{a_1},\ldots,\mysqrt{d_k}{a_k}$ where the radicands $a_i$, and exponents $d_i$ are nonnegative integers, the \emph{Radical Identity Testing} (RIT) problem is to decide whether $f(\mysqrt{d_1}{a_1},\ldots,\mysqrt{d_k}{a_k}) = 0$. The problem is easily seen to admit a {\PSPACE} upper bound by a reduction to the existential theory of reals~\cite{canny}: introduce a new formal variable for every gate of the circuit, add the equations $x_i^{d_i} - a_i = 0$ and $x_i > 0$ for every radical; RIT is now decided by checking if the resulting system of polynomial equations has a solution over the real numbers.  

Our symbolic algorithm places RIT in {\coNP}, assuming the generalised Riemann hypothesis (GRH).  The main idea behind our algorithm is that if 
$f(\mysqrt{d_1}{a_1},\ldots,\mysqrt{d_k}{a_k}) \neq 0$
then there is a polynomial-length polynomial-time checkable witness of this fact --- namely a prime $p$ and $\overline{\alpha}_1,\ldots,\overline{\alpha}_k \in \mathbb{F}_p$, satisfying $\overline{\alpha}_i ^ {d_i} \equiv a_i \pmod{p}$, such that $\overline{f}(\overline{\alpha}_1,\ldots,\overline{\alpha}_k )$ is non-zero, where $\overline{f}$ is the reduction of $f$ modulo~$p$. Crucial to our approach is the fact of \emph{joint transitivity}, which is the observation that the Galois group of the underlying real field acts jointly transitively on the roots of the various equations $x^{d_i} - a_i = 0$. This allows us to use any of the $d_i$ conjugates $\alpha_i$ of 
$\mysqrt{d_i}{a_i}$
over $\mathbb{F}_p$ in our symbolic algorithm to test whether $\overline{f}(\overline{\alpha}_1,\ldots,\overline{\alpha}_k ) = 0$. In fact, joint transitivity alone can be used in conjunction with a result of Koiran~\cite{KOIRAN1996273} to place RIT in the polynomial hierarchy assuming GRH.  Using Chebotarev's density theorem, and choosing a suitable prime, we improve this upper bound to show that RIT can be solved in {\coNP} assuming GRH.  

We next tackle a special case of RIT namely 2-RIT where the inputs to the circuit are square roots of distinct odd\footnote{
The formulation here corrects the published version~\cite{lics-version} by adding the condition that the prime be odd.} primes. The {PIT} algorithm in~\cite{chen-kao} gives a randomised numerical algorithm for the specific case of 2-RIT involving \emph{bounded} algebraic circuits\footnote{An algebraic circuit of size $s$ computing a multivariate polynomial is called bounded if the degree and bit-length of the coefficients of the polynomials computed at every gate of the circuit is bounded by a polynomial
function in~$s$.} at square roots of distinct prime numbers. However their method completely breaks down when the assumption of boundedness is dropped and the best known upper bound prior to our work was the same as that of the general RIT problem using Koiran's algorithm. We show that 2-RIT is in {\coRP} assuming GRH and in {\coNP} unconditionally. Our techniques are fundamentally different from those of ~\cite{chen-kao};  we use quadratic reciprocity and Dirichlet's theorem on the density of primes in arithmetic progressions.

\paragraph{\bf Related Work}
Balaji et al.~\cite{balaji2021cyclotomic} studied the \emph{cyclotomic identity testing}  problem, where the goal is to determine if an algebraic integer in the cyclotomic number field $\mathbb{Q}(\zeta_n)$ computed by a given algebraic circuit is zero. They show that the problem lies in the complexity class {\bpp} assuming the Generalised Riemann Hypothesis (GRH), and unconditionally in {\coNP}.
We further discuss the approach of \cite{balaji2021cyclotomic} in relation to our work in \Cref{sec:conp}.

Bl\"{o}mer~\cite{blomer98} gave a randomised algorithm to test if a bounded algebraic circuit evaluated at low degree radicals evaluates to zero. Central to such identity questions on algebraic numbers is the knowledge of the Galois group of the extension where the numbers live in. Algorithmic complexity of computing with Galois groups is investigated in ~\cite{landau-miller, landau, arvind-kurur}.

As in~\cite{blomer98}, our formulation of RIT does not permit nesting of radicals, that is, expressions such as $\sqrt{6+4\sqrt{2}}$.  Computational problems associated with nested radicals are treated in~\cite{Blomer92,Blomer97}.

\paragraph{\bf The Sum of Square Roots  problem.} Closely related to $2$-RIT is the square root sum problem where the goal is to infer the sign of a given linear combination of square roots. This is a notorious open problem in numerical analysis and computational geometry. It is known to be decidable in the \emph{Counting Hierarchy} \cite{Allender06onthe}, and the question of determining its precise computational complexity remains open since it was explicitly posed by Garey, Graham and Johnson~\cite{ggj} in 1976. Along with the related problem  of determining the sign of an integer computed by a variable-free algebraic circuit, the square root sum problem is frequently used as a tool for proving hardness and obtaining upper bounds in quantitative verification~\cite{esparza.solving, HKbudget, KWparallel, BLW}, algorithmic game theory~\cite{etessami.yannakakis, UW, CIJ}, formal language theory and logic~\cite{HKL, LOW}. We refer the interested reader to~\cite{problem33, kayal-saha} and the references therein for a discussion on the complexity status of the square root sum problem and related geometric questions.

\paragraph{\bf The Elementary Constant problem.} A related but harder identity testing problem that is a fundamental question at the intersection of transcendental number theory and model theory is the \emph{Elementary Constant problem}~\cite{richardson92}: given a complex number built from rationals using addition, multiplication and exponentiation determine if it is zero. Such numbers are called Elementary numbers, and they form an algebraically closed subfield of the complex numbers. While there is a decision procedure assuming Schanuel's conjecture~\cite{richardson97} for the elementary constant problem, the problem is not known to be decidable unconditionally, and no significant complexity lower/upper bounds are known. 


\paragraph{\bf The Compressed Word problem.} 
Arguably the earliest and most fundamental result on identity testing questions are word problems on finitely generated groups and semigroups \cite{dehn1912transformation}. There is a large body of work on the \emph{Compressed Word problem} \cite{lohrey2014compressed} which studies the computational complexity of word problems when the word is represented succinctly via a \emph{straight line program}; see \cite{konig-lohrey, balaji2021cyclotomic} for relations between such word problems and arithmetic identity testing.

